% -----------------------------------------------------------------------
% Template Tesis MKA
% 
% Diadaptasi dari https://github.com/yusufsyaifudin/template-skripsi-mipa oleh @author Yusuf Syaifudin
% @updated 09/07/2025
% 
% -----------------------------------------------------------------------

\documentclass[ugmtesis]{ugmtesis}

\input{ADDITIONAL_PACKAGES}
\titleind{JUDUL BAHASA INDONESIA}

\titleeng{JUDUL BAHASA INGGRIS}

\fullname{NAMA MAHASISWA}

\idnum{NIM MAHASISWA}

\examdate{20 Juni 2025}

\degree{\textit{Master of Computer Science (AI)}}

\yearsubmit{2025}

\program{Kecerdasan Artifisial}

\headprogram{Afiahayati, S.Kom., M.Cs., Ph.D}

\dept{Ilmu Komputer dan Elektronika}

\firstsupervisor{Nama Pembimbing Utama}

\firstexaminer{Nama Penguji Pertama}

\secondexaminer{Nama Penguji Kedua}




\usepackage{etoolbox} 
\pretocmd{\chapter}{\glsresetall}{}{}

% komen jika tidak menggunakan glosarium
\makeglossaries
% === ISTILAH ===

\newglossaryentry{finetuning}{
    name={\textit{fine-tuning}},
    description={Proses penyesuaian lanjutan model terlatih (\textit{pre-trained}) pada data tugas atau domain tertentu agar berkinerja lebih baik pada tugas tersebut}
}

\newglossaryentry{reranking}{
    name={\textit{re-ranking}},
    description={Tahap penyusunan ulang peringkat dokumen/pasase hasil pencarian tahap awal agar dokumen yang paling relevan menempati posisi teratas sebelum diberikan ke model generatif}
}

\newglossaryentry{workflow}{
    name={\textit{workflow automation}},
    description={Otomatisasi rangkaian tugas atau proses kerja yang sebelumnya dilakukan manual, pada konteks ini dijalankan oleh agen berbasis model bahasa}
}

% === SINGKATAN ===

\newacronym{ai}{AI}{\textit{Artificial Intelligence}}
\newacronym{nlp}{NLP}{\textit{Natural Language Processing}}
\newacronym{llm}{LLM}{\textit{Large Language Model}}
\newacronym{slm}{SLM}{\textit{Small Language Model}}
\newacronym{lora}{LoRA}{\textit{Low-Rank Adaptation}}
\newacronym{peft}{PEFT}{\textit{Parameter-Efficient Fine-Tuning}}
\newacronym{rag}{RAG}{\textit{Retrieval-Augmented Generation}}
\newacronym{api}{API}{\textit{Application Programming Interface}}
\newacronym{gpu}{GPU}{\textit{Graphics Processing Unit}}


\begin{document}

% Halaman judul
\titlepageind 

% Halaman Persetujuan
\approvalpage

%-----------------------------------------------------------------
% Daftar Isi
%-----------------------------------------------------------------
\newpage
\phantomsection
\addcontentsline{toc}{chapter}{\contentsname}
\tableofcontents
%-----------------------------------------------------------------
% Akhir Daftar Isi
%-----------------------------------------------------------------

%-----------------------------------------------------------------
% Daftar Tabel
%-----------------------------------------------------------------
\newpage
\phantomsection
\addcontentsline{toc}{chapter}{\listtablename}
\listoftables
%-----------------------------------------------------------------
% Akhir Daftar Tabel
%-----------------------------------------------------------------

%-----------------------------------------------------------------
% Daftar Gambar
%-----------------------------------------------------------------
\newpage
\phantomsection
\addcontentsline{toc}{chapter}{\listfigurename}
\listoffigures
%-----------------------------------------------------------------
% Akhir Daftar Gambar
%-----------------------------------------------------------------

%-----------------------------------------------------------------
% Daftar Algoritma
%-----------------------------------------------------------------
% dikomen jika tidak menggunakan algoritma
\newpage
\phantomsection
\renewcommand{\listalgorithmname}{DAFTAR ALGORITMA}
\addcontentsline{toc}{chapter}{DAFTAR ALGORITMA}
\listofalgorithms % Menampilkan daftar algoritma
%-----------------------------------------------------------------
% Akhir Daftar Algoritma
%-----------------------------------------------------------------


%-----------------------------------------------------------------
% Daftar Kode Program
%-----------------------------------------------------------------
\newpage
\phantomsection
% \renewcommand{\thechapter}{\arabic{chapter}}
\renewcommand{\thelstlisting}{\thechapter.\arabic{lstlisting}}
\renewcommand\lstlistingname{Kode Program}
\renewcommand\lstlistlistingname{DAFTAR KODE PROGRAM}
\addcontentsline{toc}{chapter}{DAFTAR KODE PROGRAM}
\lstlistoflistings


%-----------------------------------------------------------------
% Akhir Daftar Kode Program
%-----------------------------------------------------------------

%-----------------------------------------------------------------
% Daftar Istilah
%-----------------------------------------------------------------
\newpage
\phantomsection
\addcontentsline{toc}{chapter}{DAFTAR ISTILAH}
\printglossary[type=main, title=DAFTAR ISTILAH]
%-----------------------------------------------------------------
% Akhir Daftar Istilah
%-----------------------------------------------------------------

%-----------------------------------------------------------------
% Daftar Singkatan
%-----------------------------------------------------------------
\newpage
\phantomsection
\addcontentsline{toc}{chapter}{DAFTAR SINGKATAN}
\printglossary[type=\acronymtype, title=DAFTAR SINGKATAN]
%-----------------------------------------------------------------
% Akhir Singkatan
%-----------------------------------------------------------------

%-----------------------------------------------------------------
%Disini awal masukan Intisari
%-----------------------------------------------------------------
\begin{abstractind}
	Pemanfaatan \gls{ai} \textit{agent} untuk otomatisasi alur kerja (\gls{workflow}) berkembang pesat, namun sebagian besar masih bergantung pada \gls{llm} berukuran sangat besar yang berjalan di layanan awan berbayar sehingga menimbulkan persoalan biaya, latensi, serta privasi data. \gls{slm} yang dapat dijalankan secara lokal menjadi alternatif yang menarik, tetapi akurasinya pada tugas penalaran dan pemanggilan \textit{tool} masih tertinggal dibanding model besar. Penelitian ini mengusulkan optimasi akurasi \gls{slm} lokal untuk \gls{ai} \textit{agent workflow automation} melalui kombinasi dua pendekatan: (1) \gls{finetuning} hemat parameter dengan \gls{lora} untuk adaptasi domain, dan (2) \gls{rag} yang diperkuat mekanisme \gls{reranking} agar konteks paling relevan diberikan kepada model yang memiliki keterbatasan panjang konteks. Penelitian akan mengevaluasi kontribusi tiap komponen dan kombinasinya melalui studi ablasi terhadap metrik akurasi tugas, kualitas retrieval, serta efisiensi komputasi. Kontribusi yang diharapkan adalah rancangan \textit{pipeline} \gls{slm} lokal yang lebih akurat namun tetap efisien untuk diterapkan pada otomatisasi alur kerja berbasis agen.

\vspace{0.5\baselineskip}
\noindent\textbf{Kata Kunci}: \textit{small language model}, \textit{LoRA}, \textit{retrieval-augmented generation}, \textit{re-ranking}, \textit{AI agent}

\end{abstractind}
%-----------------------------------------------------------------
%Disini akhir masukan Intisari
%-----------------------------------------------------------------

%-----------------------------------------------------------------
% Konten Tesis
%-----------------------------------------------------------------
\chapter{PENDAHULUAN}
\label{PENDAHULUAN}
\section{Latar Belakang}
Kemajuan \gls{llm} dalam beberapa tahun terakhir telah mendorong perubahan besar pada cara sistem perangkat lunak dibangun, termasuk munculnya paradigma \gls{ai} \textit{agent}, yaitu sistem yang memanfaatkan model bahasa untuk menalar, merencanakan, dan memanggil \textit{tool} eksternal guna menyelesaikan tugas secara otomatis \citep{minaee2024llmsurvey, wang2024agentsurvey}. Kemampuan model bahasa untuk memutuskan kapan dan bagaimana menggunakan \gls{api} atau alat bantu eksternal \citep{schick2023toolformer}, serta menalar secara bertahap untuk memecahkan masalah kompleks \citep{yao2023tot}, menjadikan agen berbasis model bahasa sebagai fondasi yang menjanjikan bagi \gls{workflow} pada berbagai domain \citep{park2023generativeagents}.

Namun, sebagian besar agen otomatisasi alur kerja saat ini bertumpu pada \gls{llm} berukuran sangat besar yang umumnya diakses sebagai layanan awan berbayar. Ketergantungan ini menimbulkan sejumlah persoalan praktis, antara lain biaya pemanggilan yang tinggi pada penggunaan berskala besar, latensi jaringan, serta risiko privasi dan kepatuhan ketika data sensitif organisasi harus dikirim ke pihak ketiga. Persoalan tersebut mendorong kebutuhan akan model bahasa yang dapat dijalankan secara lokal (\textit{on-premise}) sehingga organisasi memiliki kendali penuh atas data dan biaya operasionalnya.

\gls{slm}, yaitu model bahasa dengan jumlah parameter relatif kecil sehingga dapat dijalankan pada perangkat keras terbatas, hadir sebagai alternatif yang menarik untuk kebutuhan tersebut \citep{wang2025slmsurvey}. Sejumlah \gls{slm} terbuka terbaru menunjukkan kemampuan yang menjanjikan meskipun berukuran kecil, seperti Phi-3 yang dirancang agar dapat berjalan di perangkat ponsel \citep{abdin2024phi3}, TinyLlama \citep{zhang2024tinyllama}, serta keluarga model Qwen2.5 dengan berbagai ukuran \citep{qwen2024qwen25}. Walaupun demikian, pada tugas yang menuntut penalaran, kepatuhan terhadap instruksi, dan pemanggilan \textit{tool} secara tepat sebagaimana dibutuhkan dalam alur kerja agen, akurasi \gls{slm} masih tertinggal dibandingkan model berukuran besar \citep{wang2025slmsurvey}. Oleh karena itu, dibutuhkan strategi untuk meningkatkan akurasi \gls{slm} lokal agar layak digunakan pada otomatisasi alur kerja berbasis agen.

Salah satu strategi yang efektif untuk meningkatkan kinerja model pada domain spesifik adalah \gls{finetuning}. Namun, \gls{finetuning} penuh terhadap seluruh parameter model semakin tidak efisien seiring bertambahnya ukuran model \citep{ding2023delta}. Pendekatan \gls{peft}, khususnya \gls{lora}, menjawab persoalan ini dengan membekukan bobot model terlatih dan hanya melatih sejumlah kecil matriks berperingkat rendah, sehingga jumlah parameter yang dilatih berkurang drastis tanpa menambah latensi inferensi \citep{hu2022lora}. Varian seperti QLoRA bahkan memungkinkan \gls{finetuning} model pada satu \gls{gpu} dengan memori terbatas melalui kuantisasi \citep{dettmers2023qlora}, sehingga \gls{lora} sangat sesuai untuk konteks pengembangan \gls{slm} lokal dengan sumber daya komputasi yang hemat \citep{mao2024lorasurvey}.

Strategi kedua untuk meningkatkan akurasi tanpa memperbesar model adalah \gls{rag}, yang memadukan model generatif dengan memori non-parametrik berupa basis pengetahuan eksternal yang dapat diakses saat inferensi \citep{lewis2020rag}. \gls{rag} terbukti mengurangi halusinasi dan memungkinkan pembaruan pengetahuan tanpa melatih ulang model \citep{gao2023ragsurvey}. Akan tetapi, efektivitas \gls{rag} sangat bergantung pada kualitas dokumen yang diambil. Model bahasa, terutama yang berukuran kecil dengan jendela konteks terbatas, terbukti tidak robust dalam memanfaatkan informasi yang posisinya berada di tengah konteks panjang, suatu fenomena yang dikenal sebagai \textit{lost in the middle} \citep{liu2024lostmiddle}. Hal ini menyiratkan bahwa sekadar menambahkan banyak dokumen hasil pencarian tahap awal belum tentu meningkatkan akurasi; diperlukan tahap \gls{reranking} untuk memastikan dokumen yang paling relevan menempati posisi teratas. Penelitian terkini menunjukkan bahwa \gls{reranking}, baik dengan pendekatan terlatih \citep{ren2021rocketqav2} maupun berbasis model bahasa \citep{sun2023rankgpt}, dapat meningkatkan relevansi konteks secara signifikan.

Meskipun \gls{lora} dan \gls{rag} dengan \gls{reranking} masing-masing telah banyak diteliti, kombinasi keduanya untuk meningkatkan akurasi \gls{slm} lokal pada konteks \gls{ai} \textit{agent workflow automation} masih belum dieksplorasi secara sistematis. Di satu sisi, pemanfaatan \gls{slm} untuk alur kerja agen masih terbatas: \citet{erdogan2024tinyagent} menunjukkan agen \gls{slm} yang hanya melakukan \textit{function-calling} tanpa \gls{rag}, sedangkan kerangka orkestrasi mutakhir seperti \citet{niu2025flow} dan \citet{han2026legomem} berfokus pada sistem multi-agen berbasis \gls{llm} besar. Di sisi lain, peningkatan \gls{rag} melalui \gls{reranking} umumnya dilakukan terpisah dari \gls{finetuning} model \citep{yu2024rankrag, saha2024quimrag}. Akibatnya, belum jelas seberapa besar kontribusi masing-masing komponen, bagaimana interaksinya, serta konfigurasi \gls{reranking} mana yang memberikan keseimbangan terbaik antara akurasi dan efisiensi pada perangkat dengan sumber daya terbatas. Penelitian ini diajukan untuk menjawab kekosongan tersebut, yaitu merancang dan mengevaluasi \textit{pipeline} yang memadukan \gls{finetuning} \gls{lora} dan \gls{rag} dengan \gls{reranking} untuk mengoptimalkan akurasi \gls{slm} lokal pada tugas otomatisasi alur kerja berbasis agen.

% TODO (Aditya): pertajam latar belakang dengan domain/kasus uji spesifik yang akan dipakai
% (mis. otomatisasi tiket layanan, asisten basis pengetahuan internal, dsb.) bila sudah ditetapkan.

\section{Rumusan Masalah}
Berdasarkan latar belakang di atas, permasalahan utama penelitian ini dirumuskan sebagai berikut: \textit{bagaimana mengoptimalkan akurasi \gls{slm} lokal menggunakan \gls{finetuning} \gls{lora} dan \gls{rag} ber-\gls{reranking} pada \gls{ai} \textit{agent} untuk otomatisasi alur kerja (\gls{workflow})?} Permasalahan utama tersebut diuraikan menjadi pertanyaan penelitian berikut:
\begin{enumerate}
    \item Seberapa besar peningkatan akurasi \gls{slm} lokal pada tugas \gls{ai} \textit{agent workflow automation} ketika diterapkan \gls{finetuning} \gls{lora} dibandingkan model dasar (\textit{baseline}) tanpa \gls{finetuning}?
    \item Bagaimana pengaruh penerapan \gls{rag} dengan \gls{reranking} terhadap akurasi \gls{slm} lokal, dan bagaimana kinerja kombinasi \gls{lora} dengan \gls{rag} ber-\gls{reranking} dibandingkan penerapan masing-masing komponen secara terpisah?
    \item Konfigurasi \gls{reranking} manakah yang memberikan keseimbangan terbaik antara akurasi dan efisiensi komputasi agar tetap layak dijalankan pada lingkungan lokal dengan sumber daya terbatas?
\end{enumerate}

\section{Batasan Masalah}
Agar penelitian terarah, ditetapkan batasan sebagai berikut:
\begin{enumerate}
    \item Model yang dioptimasi adalah \gls{slm} terbuka berukuran kecil yang dapat dijalankan secara lokal; model bahasa berukuran sangat besar dan layanan berbayar hanya digunakan sebagai pembanding bila diperlukan.
    \item Teknik \gls{finetuning} yang dikaji difokuskan pada \gls{lora} dan variannya, bukan \gls{finetuning} penuh.
    \item Peningkatan akurasi dievaluasi pada tugas otomatisasi alur kerja berbasis agen pada domain dan dataset yang ditetapkan dalam penelitian ini. % TODO: sebutkan dataset/domain final
    \item Penelitian berfokus pada akurasi, kualitas \textit{retrieval}, dan efisiensi komputasi, serta tidak mencakup aspek antarmuka pengguna maupun penerapan produksi berskala penuh.
\end{enumerate}

\section{Tujuan Penelitian}
Tujuan yang ingin dicapai pada penelitian ini adalah:
\begin{enumerate}
    \item Mengukur dan menganalisis peningkatan akurasi \gls{slm} lokal akibat penerapan \gls{finetuning} \gls{lora} pada tugas otomatisasi alur kerja berbasis agen.
    \item Menganalisis pengaruh \gls{rag} dengan \gls{reranking} serta kombinasinya dengan \gls{lora} terhadap akurasi melalui studi ablasi.
    \item Menentukan konfigurasi \gls{reranking} yang memberikan keseimbangan optimal antara akurasi dan efisiensi komputasi untuk penerapan lokal.
\end{enumerate}

\section{Manfaat Penelitian}
Penelitian ini diharapkan memberi manfaat sebagai berikut:
\begin{enumerate}
    \item \textbf{Manfaat teoretis:} memberikan bukti empiris mengenai kontribusi dan interaksi \gls{lora} serta \gls{rag} ber-\gls{reranking} dalam meningkatkan akurasi \gls{slm}, yang memperkaya khazanah penelitian \gls{nlp} hemat sumber daya.
    \item \textbf{Manfaat praktis:} menghasilkan rancangan \textit{pipeline} \gls{slm} lokal yang lebih akurat dan efisien sehingga dapat menjadi acuan pengembangan \gls{ai} \textit{agent} yang menjaga privasi data dan menekan biaya operasional.
    \item \textbf{Manfaat metodologis:} menyediakan prosedur evaluasi yang reprodusibel untuk membandingkan strategi optimasi \gls{slm} pada tugas otomatisasi alur kerja.
\end{enumerate}


\chapter{TINJAUAN PUSTAKA}
\label{TINJAUAN PUSTAKA}
Bab ini menyajikan tinjauan sistematis atas penelitian terdahulu yang relevan dengan optimasi akurasi \gls{slm}, dikelompokkan menurut tema utama, lalu diakhiri dengan posisi penelitian ini terhadap karya-karya tersebut. Sitasi menggunakan format nama--tahun.

\section{Small Language Model}
Tren penelitian model bahasa bergeser dari sekadar memperbesar ukuran menuju efisiensi, melahirkan kelas \gls{slm} yang menyeimbangkan kemampuan dan biaya komputasi \citep{wang2025slmsurvey}. \citet{abdin2024phi3} menunjukkan bahwa model 3,8 miliar parameter (Phi-3-mini) yang dilatih dengan kurasi data berkualitas dapat menyaingi model yang jauh lebih besar dan tetap mampu berjalan di perangkat ponsel. Pendekatan lain mencapai efisiensi melalui pra-pelatihan ringkas seperti TinyLlama \citep{zhang2024tinyllama}, penyediaan beragam ukuran model terbuka seperti Qwen2.5 \citep{qwen2024qwen25}, maupun pemangkasan terstruktur dari model besar seperti Sheared-LLaMA \citep{xia2023sheared}. Karya-karya ini menegaskan kelayakan \gls{slm} untuk penerapan lokal, sekaligus menyoroti bahwa kesenjangan akurasi terhadap model besar masih menjadi tantangan terbuka.

\section{Parameter-Efficient Fine-Tuning dan LoRA}
\citet{ding2023delta} memberikan tinjauan menyeluruh atas \gls{peft} dan menunjukkan bahwa pengoptimalan sebagian kecil parameter dapat mendekati kinerja \gls{finetuning} penuh. Metode \gls{lora} \citep{hu2022lora} menjadi tonggak penting dengan menyuntikkan matriks dekomposisi berperingkat rendah pada arsitektur \textit{transformer}. Pengembangan lanjutan mencakup QLoRA yang menggabungkan kuantisasi 4-bit untuk efisiensi memori \citep{dettmers2023qlora}, DyLoRA yang menghindari pencarian peringkat secara manual \citep{valipour2023dylora}, serta kerangka adapter terpadu untuk \gls{slm} \citep{hu2023llmadapters}. Ketersediaan perkakas seperti LlamaFactory \citep{zheng2024llamafactory} mempermudah penerapan teknik-teknik tersebut secara praktis. Tinjauan khusus oleh \citet{mao2024lorasurvey} memetakan berbagai varian \gls{lora} beserta arah pengembangannya.

\section{Retrieval-Augmented Generation}
\citet{lewis2020rag} memperkenalkan \gls{rag} yang memadukan memori parametrik dan non-parametrik untuk tugas padat pengetahuan. Tinjauan oleh \citet{gao2023ragsurvey} mengklasifikasikan perkembangan \gls{rag} dari \textit{naive}, \textit{advanced}, hingga \textit{modular}. \citet{ram2023incontext} menunjukkan bahwa pendekatan In-Context RALM mampu meningkatkan kinerja model bahasa tanpa mengubah arsitekturnya, sementara \citet{jiang2023activerag} mengusulkan pengambilan dokumen secara aktif selama proses generasi. Untuk pengukuran, \citet{chen2024ragbench} dan \citet{es2024ragas} menyediakan tolok ukur serta metrik evaluasi otomatis bagi sistem \gls{rag}.

\section{Re-ranking pada Sistem Retrieval}
Kualitas konteks yang diberikan kepada model generatif sangat menentukan akurasi keluaran. \citet{liu2024lostmiddle} membuktikan bahwa model bahasa kurang efektif memanfaatkan informasi yang terletak di tengah konteks panjang, sehingga urutan dokumen menjadi krusial. \citet{ren2021rocketqav2} mengusulkan pelatihan gabungan \textit{retriever} dan \textit{re-ranker} secara terpadu. Perkembangan terbaru memanfaatkan model bahasa sebagai \textit{re-ranker}: \citet{sun2023rankgpt} menunjukkan bahwa model bahasa dapat menjadi agen pemeringkat yang kompetitif dan kemampuannya dapat didistilasi ke model kecil, sedangkan \citet{ma2023lrl}, \citet{pradeep2023rankvicuna}, dan \citet{pradeep2023rankzephyr} mengembangkan \textit{re-ranker} \textit{listwise} berbasis model terbuka.

\section{AI Agent dan Workflow Automation}
\citet{wang2024agentsurvey} menyusun tinjauan komprehensif arsitektur agen otonom berbasis model bahasa. Kemampuan inti agen mencakup penggunaan \textit{tool} eksternal \citep{schick2023toolformer}, penalaran bertahap yang lebih terstruktur \citep{yao2023tot}, serta perencanaan dan memori sebagaimana ditunjukkan pada agen generatif \citep{park2023generativeagents}. Karya-karya ini menjadi landasan bagi penerapan model bahasa pada otomatisasi alur kerja.

\section{Posisi Penelitian}
Tinjauan di atas menunjukkan bahwa \gls{lora}, \gls{rag}, dan \gls{reranking} masing-masing efektif meningkatkan kinerja model bahasa, namun umumnya diteliti secara terpisah dan kerap pada model berukuran besar. Belum banyak penelitian yang secara sistematis memadukan \gls{finetuning} \gls{lora} dengan \gls{rag} ber-\gls{reranking} khusus untuk meningkatkan akurasi \gls{slm} lokal pada konteks \gls{ai} \textit{agent workflow automation}, sekaligus menganalisis kontribusi tiap komponen dan keseimbangannya dengan efisiensi komputasi. Kekosongan inilah yang menjadi fokus penelitian ini. Ringkasan sebagian penelitian terdahulu yang menjadi acuan utama disajikan pada Tabel \ref{tab:tinjauan-pustaka}.

\begin{longtable}{p{0.16\linewidth} p{0.24\linewidth} p{0.26\linewidth} p{0.24\linewidth}}
    \caption{Ringkasan Penelitian Terdahulu yang Menjadi Acuan Utama}
    \label{tab:tinjauan-pustaka} \\
    \hline
    \textbf{Penelitian} & \textbf{Tujuan} & \textbf{Metode} & \textbf{Hasil Utama} \\
    \hline
    \endfirsthead

    \hline
    \textbf{Penelitian} & \textbf{Tujuan} & \textbf{Metode} & \textbf{Hasil Utama} \\
    \hline
    \endhead

    \hline
    \multicolumn{4}{r}{\textit{Lanjutan di halaman berikutnya...}} \\
    \endfoot

    \hline
    \endlastfoot

    \citet{hu2022lora} &
    Adaptasi LLM secara hemat parameter &
    Menyuntikkan matriks berperingkat rendah, bobot dasar dibekukan &
    Parameter terlatih turun drastis, setara/lebih baik dari \textit{fine-tuning} penuh tanpa latensi tambahan \\
    \hline

    \citet{dettmers2023qlora} &
    \textit{Fine-tuning} hemat memori &
    Kuantisasi 4-bit (NF4) dipadu LoRA dan \textit{paged optimizer} &
    Mampu melatih model besar pada satu GPU dengan kualitas mendekati 16-bit \\
    \hline

    \citet{lewis2020rag} &
    Tugas NLP padat pengetahuan &
    Gabungan model seq2seq pra-latih dengan \textit{retriever} \textit{dense} &
    Capaian SOTA pada beberapa tugas QA domain terbuka, keluaran lebih faktual \\
    \hline

    \citet{ram2023incontext} &
    Meningkatkan LM tanpa ubah arsitektur &
    In-Context RALM dengan spesialisasi mekanisme \textit{ranking} &
    Peningkatan kinerja konsisten lintas ukuran model \\
    \hline

    \citet{sun2023rankgpt} &
    \textit{Re-ranking} relevansi dokumen &
    Model bahasa sebagai agen pemeringkat \textit{listwise} dan distilasi permutasi &
    Kompetitif/melebihi metode \textit{supervised}; model distilasi kecil mengungguli model \textit{supervised} lebih besar \\
    \hline

    \citet{abdin2024phi3} &
    SLM berkemampuan tinggi di perangkat lokal &
    Model 3,8 M parameter dengan kurasi data berkualitas &
    Setara model jauh lebih besar pada sejumlah tolok ukur, dapat berjalan di ponsel \\
    \hline

\end{longtable}


\chapter{DASAR TEORI}
\label{DASAR TEORI}
Bab ini menguraikan landasan teori yang menjadi dasar penelitian, meliputi arsitektur model bahasa, \gls{slm}, \gls{peft} dengan \gls{lora}, \gls{rag}, mekanisme \gls{reranking}, serta konsep \gls{ai} \textit{agent}.

\section{Model Bahasa dan Arsitektur \textit{Transformer}}
Model bahasa modern dibangun di atas arsitektur \textit{transformer} yang mengandalkan mekanisme \textit{self-attention} untuk menangkap keterkaitan antar-token dalam suatu urutan \citep{minaee2024llmsurvey}. Inti dari mekanisme tersebut adalah operasi \textit{attention} yang dirumuskan pada Persamaan \ref{eq:attention}, dengan $Q$, $K$, dan $V$ berturut-turut adalah matriks \textit{query}, \textit{key}, dan \textit{value}, serta $d_k$ adalah dimensi \textit{key}.

\begin{equation}
    \text{Attention}(Q, K, V) = \text{softmax}\left( \frac{QK^\top}{\sqrt{d_k}} \right) V
    \label{eq:attention}
\end{equation}

Model bahasa berskala besar dilatih pada korpus teks masif sehingga mampu menyimpan pengetahuan dalam parameternya dan menunjukkan kemampuan generalisasi yang kuat \citep{minaee2024llmsurvey}.

\section{Small Language Model}
\gls{slm} adalah model bahasa dengan jumlah parameter relatif kecil (umumnya pada kisaran ratusan juta hingga beberapa miliar) sehingga dapat dijalankan pada perangkat keras terbatas atau secara lokal \citep{wang2025slmsurvey}. Efisiensi \gls{slm} dicapai melalui kombinasi arsitektur yang ringkas, kurasi data pra-pelatihan yang berkualitas, serta teknik kompresi. Keunggulan utama \gls{slm} adalah biaya inferensi yang rendah, latensi yang kecil, dan kemampuan berjalan tanpa ketergantungan pada layanan awan, dengan konsekuensi kapasitas penalaran yang lebih terbatas dibanding model besar \citep{wang2025slmsurvey}.

\section{Parameter-Efficient Fine-Tuning dan LoRA}
\gls{peft} adalah keluarga teknik yang mengadaptasi model terlatih dengan hanya memperbarui sebagian kecil parameter, sementara sebagian besar bobot dibekukan \citep{ding2023delta}. \gls{lora} \citep{hu2022lora} merupakan salah satu metode \gls{peft} paling banyak digunakan. Diberikan bobot terlatih $W_0 \in \mathbb{R}^{d \times k}$, \gls{lora} memodelkan pembaruan bobot $\Delta W$ sebagai hasil perkalian dua matriks berperingkat rendah, sebagaimana Persamaan \ref{eq:lora}, dengan $B \in \mathbb{R}^{d \times r}$, $A \in \mathbb{R}^{r \times k}$, dan peringkat $r \ll \min(d,k)$.

\begin{equation}
    h = W_0 x + \Delta W x = W_0 x + B A x
    \label{eq:lora}
\end{equation}

Karena hanya $A$ dan $B$ yang dilatih, jumlah parameter yang diperbarui jauh lebih kecil daripada \textit{fine-tuning} penuh, tanpa menambah latensi inferensi karena $B A$ dapat digabungkan kembali ke $W_0$ saat penerapan \citep{hu2022lora}. Varian QLoRA menambahkan kuantisasi 4-bit pada bobot dasar untuk menekan kebutuhan memori \citep{dettmers2023qlora}.

\section{Retrieval-Augmented Generation}
\gls{rag} memadukan kemampuan generatif model bahasa (memori parametrik) dengan basis pengetahuan eksternal yang dapat diambil saat inferensi (memori non-parametrik) \citep{lewis2020rag}. Secara umum, alur \gls{rag} terdiri atas tiga tahap: (1) \textit{retrieval}, yaitu pengambilan dokumen relevan dari basis pengetahuan menggunakan representasi vektor; (2) \textit{augmentation}, yaitu penggabungan dokumen terambil ke dalam \textit{prompt}; dan (3) \textit{generation}, yaitu penghasilan jawaban oleh model bahasa berdasarkan konteks tersebut \citep{gao2023ragsurvey}. Pendekatan ini memungkinkan pembaruan pengetahuan tanpa pelatihan ulang dan mengurangi halusinasi.

\section{Re-ranking}
\gls{reranking} adalah tahap menyusun ulang peringkat dokumen hasil \textit{retrieval} awal agar dokumen paling relevan menempati posisi teratas sebelum diberikan ke model generatif. Tahap ini penting karena model bahasa, khususnya yang berukuran kecil dengan konteks terbatas, terbukti tidak robust memanfaatkan informasi yang terletak di tengah konteks panjang \citep{liu2024lostmiddle}. \gls{reranking} dapat dilakukan dengan model terlatih khusus \citep{ren2021rocketqav2} maupun dengan memanfaatkan model bahasa sebagai pemeringkat \textit{listwise} \citep{sun2023rankgpt}.

\section{AI Agent}
\gls{ai} \textit{agent} berbasis model bahasa adalah sistem yang menggunakan model bahasa sebagai \textit{controller} untuk menalar, merencanakan, dan berinteraksi dengan lingkungan melalui \textit{tool} eksternal guna menyelesaikan tugas \citep{wang2024agentsurvey}. Komponen pentingnya mencakup kemampuan pemanggilan \textit{tool}/\gls{api} \citep{schick2023toolformer} dan penalaran bertahap \citep{yao2023tot}. Dalam konteks \gls{workflow}, agen memetakan instruksi pengguna menjadi rangkaian tindakan terotomasi, sehingga akurasi model dalam memahami instruksi dan memilih tindakan yang tepat menjadi faktor kritis.


\chapter{METODOLOGI PENELITIAN}
\label{ANALISIS DAN PERANCANGAN SISTEM}
Bab ini menjelaskan metodologi penelitian, mencakup deskripsi umum, akuisisi data, rancangan model yang diusulkan, serta rancangan pengujian. Penelitian ini bersifat eksperimental kuantitatif dengan pendekatan pengembangan dan evaluasi sistem.

\section{Deskripsi Umum Penelitian}
Penelitian ini bertujuan mengoptimalkan akurasi \gls{slm} lokal untuk \gls{ai} \textit{agent workflow automation} melalui kombinasi \gls{finetuning} \gls{lora} dan \gls{rag} ber-\gls{reranking}. Alur penelitian terdiri atas lima tahap: (1) penyiapan data dan model dasar; (2) penerapan \gls{finetuning} \gls{lora}; (3) pembangunan \textit{pipeline} \gls{rag} dengan \gls{reranking}; (4) integrasi dan eksperimen ablasi; serta (5) evaluasi dan analisis. Variabel bebas penelitian adalah strategi optimasi (\gls{lora}, \gls{rag}, \gls{reranking}, dan kombinasinya), sedangkan variabel terikat adalah akurasi tugas, kualitas \textit{retrieval}, dan efisiensi komputasi.

\section{Akuisisi Data}
Data yang digunakan meliputi tiga komponen:
\begin{enumerate}
    \item \textbf{Model dasar (\textit{base} \gls{slm}):} \gls{slm} terbuka yang dapat dijalankan lokal, dengan kandidat antara lain Phi-3-mini \citep{abdin2024phi3}, Qwen2.5 berukuran kecil \citep{qwen2024qwen25}, atau TinyLlama \citep{zhang2024tinyllama}. % TODO (Aditya): tetapkan model final beserta alasannya.
    \item \textbf{Data \gls{finetuning}:} kumpulan data instruksi untuk tugas otomatisasi alur kerja agen (mis. pemanggilan \textit{tool}, penalaran prosedural, dan menjawab berbasis domain). % TODO: tetapkan sumber/format dataset dan jumlah sampel.
    \item \textbf{Basis pengetahuan \gls{rag}:} korpus dokumen domain yang diindeks sebagai basis \textit{retrieval}. % TODO: tetapkan korpus domain.
\end{enumerate}
Data dibagi menjadi himpunan latih, validasi, dan uji dengan proporsi yang ditetapkan secara konsisten, dan seluruh praproses (pembersihan, tokenisasi, penyusunan \textit{prompt}) didokumentasikan agar reprodusibel.

\section{Rancangan Model}
Sistem yang diusulkan terdiri atas dua jalur optimasi yang saling melengkapi. Pertama, model dasar diadaptasi ke domain melalui \gls{finetuning} \gls{lora} \citep{hu2022lora}, bila perlu dengan kuantisasi mengikuti skema QLoRA untuk menghemat memori \citep{dettmers2023qlora}; implementasi pelatihan memanfaatkan perkakas \gls{finetuning} yang telah mapan \citep{zheng2024llamafactory}. Kedua, pada tahap inferensi, model di-augmentasi dengan \gls{rag} \citep{lewis2020rag, gao2023ragsurvey} yang dilengkapi tahap \gls{reranking} untuk memastikan konteks paling relevan diberikan kepada model, mengingat keterbatasan pemanfaatan konteks panjang pada model kecil \citep{liu2024lostmiddle}. Mekanisme \gls{reranking} yang dibandingkan mencakup pendekatan \textit{re-ranker} terlatih \citep{ren2021rocketqav2} dan pendekatan berbasis model bahasa \textit{listwise} \citep{sun2023rankgpt}. Garis besar alur inferensi diuraikan pada \textit{pseudocode} \ref{alg:pipeline}.

\begin{algorithm}[H]
\caption{Alur Inferensi SLM dengan LoRA dan Re-ranked RAG}
\label{alg:pipeline}
\begin{algorithmic}[1]
\Procedure{JawabAgen}{kueri $q$, basis pengetahuan $\mathcal{D}$, model $M_{\text{LoRA}}$}
    \State $C \gets \text{Retrieve}(q, \mathcal{D}, k)$ \Comment{ambil $k$ dokumen kandidat}
    \State $C' \gets \text{ReRank}(q, C, k')$ \Comment{susun ulang, ambil $k' \le k$ teratas}
    \State $p \gets \text{SusunPrompt}(q, C')$ \Comment{gabungkan konteks terpilih}
    \State $a \gets M_{\text{LoRA}}(p)$ \Comment{hasilkan jawaban/tindakan agen}
    \State \Return $a$
\EndProcedure
\end{algorithmic}
\end{algorithm}

Konfigurasi yang dieksplorasi mencakup peringkat \gls{lora} ($r$), faktor penskalaan, dan modul sasaran; jumlah dokumen yang diambil ($k$) dan disisakan setelah \gls{reranking} ($k'$); serta jenis \textit{re-ranker}. % TODO: tetapkan rentang nilai hyperparameter final.

\section{Rancangan Pengujian}
Untuk menjawab rumusan masalah, dilakukan studi ablasi yang membandingkan beberapa konfigurasi: (K1) model dasar tanpa intervensi sebagai \textit{baseline}; (K2) model dasar + \gls{lora}; (K3) model dasar + \gls{rag} tanpa \gls{reranking}; (K4) model dasar + \gls{rag} + \gls{reranking}; dan (K5) model dasar + \gls{lora} + \gls{rag} + \gls{reranking}. Perbandingan antar-konfigurasi memungkinkan pengukuran kontribusi tiap komponen dan interaksinya.

Metrik evaluasi yang digunakan dikelompokkan menjadi tiga, sebagaimana Tabel \ref{tab:metrik}.

\begin{table}[!h]
    \centering
    \caption{Metrik Evaluasi}
    \label{tab:metrik}
    \begin{tabular}{p{0.26\linewidth} p{0.62\linewidth}}
        \toprule
        \textbf{Aspek} & \textbf{Metrik} \\
        \midrule
        Akurasi tugas & \textit{Exact Match}, F1, dan/atau akurasi pemanggilan \textit{tool} sesuai tugas agen \\
        Kualitas \textit{retrieval}/jawaban & nDCG/\textit{recall} untuk \textit{retrieval}; metrik berbasis \gls{rag} seperti \textit{faithfulness} dan \textit{answer relevance} \citep{es2024ragas, chen2024ragbench} \\
        Efisiensi & Latensi inferensi, penggunaan memori \gls{gpu}, dan jumlah parameter terlatih \\
        \bottomrule
    \end{tabular}
\end{table}

Eksperimen dijalankan pada lingkungan perangkat keras yang ditetapkan, dengan \textit{seed} acak yang dikunci, konfigurasi yang disimpan sebagai berkas, dan kode yang ditandai versinya untuk menjamin reprodusibilitas. % TODO: cantumkan spesifikasi perangkat keras (GPU) yang dipakai.


\chapter{JADWAL PENELITIAN}
\label{IMPLEMENTASI SISTEM}
Bagian ini menyajikan rencana jadwal penelitian beserta tahapan dan estimasi waktu pelaksanaan, mengikuti rencana yang telah dipresentasikan dan disetujui pembimbing. Jadwal dapat disesuaikan mengikuti kalender akademik. Matriks kegiatan ditunjukkan pada Tabel \ref{tab:jadwal_penelitian}.

\begin{table}[!h]
    \centering
    \caption{Jadwal Kegiatan Penelitian}
    \label{tab:jadwal_penelitian}
    \resizebox{\linewidth}{!}{
    \begin{tabular}{clcccccccccc}
    \toprule
        No & Kegiatan & Mei & Jun & Jul & Agu & Sep & Okt & Nov & Des & Jan & Feb \\
    \midrule
        1 & Identifikasi masalah \& studi literatur & X & X & & & & & & & & \\
        2 & Penyusunan proposal tesis & & X & X & X & & & & & & \\
        3 & Pengumpulan \& penyusunan dataset & & & & X & X & X & & & & \\
        4 & \textit{Fine-tuning} \textit{Local} SLM (LoRA) & & & & & & X & X & & & \\
        5 & Implementasi RAG \& \textit{re-ranking} & & & & & & & X & X & X & \\
        6 & Pengujian \& evaluasi sistem & & & & & & & & X & X & \\
        7 & Analisis hasil penelitian & & & & & & & & & X & X \\
        8 & Penulisan tesis \& revisi & & & & & & & & X & X & X \\
    \bottomrule
    \end{tabular}
    }
\end{table}

Keterangan: tanda ``X'' menunjukkan periode pelaksanaan tiap kegiatan. % TODO (Aditya): sesuaikan tahun dan bulan mengikuti kalender akademik serta arahan pembimbing.



%-----------------------------------------------------------------
% Awal Daftar Pustaka
%-----------------------------------------------------------------
\addcontentsline{toc}{chapter}{DAFTAR PUSTAKA}
\bibliography{referensi}
%-----------------------------------------------------------------
% Akhir Daftar Pustaka
%-----------------------------------------------------------------

%-----------------------------------------------------------------
% Awal lampiran
%-----------------------------------------------------------------
% dibuka komennya kalau ada lampirannya
% \appendix
% \chapter{BERKAS JSON UNTUK MODEL SISTEM PENGENALAN ENTITAS BERNAMA}
% \label{BERKAS JSON UNTUK MODEL SISTEM PENGENALAN ENTITAS BERNAMA}
% \input{BAB_TESIS/BAB9_LAMPIRAN/1_BERKAS_MODEL_NER}

%-----------------------------------------------------------------
% Akhir lampiran
%-----------------------------------------------------------------

\end{document}